% GENERATED FILE. Edit canonical lessons, then run npm run generate:books.
% canonical-source-hash: fnv1a64:a04c6211888b986a
% canonical-lessons: HI-C475-caritra, HI-C475-acran, HI-C475-mizaj, HI-C475-gun, HI-C475-pravrtti, HI-R475-review-468-475

\chapter{Character and Conduct}
\label{ch:hi-character-and-conduct}
\hlchaptermodality{475}

\begin{chapteropening}
\textbf{By the end of this chapter:} \emph{I can describe someone's character and moods.}

Complete the last lesson of chapter 475: I can describe someone's character and moods.
\end{chapteropening}

\section[\texorpdfstring{\hi{चरित्र} (caritra)}{चरित्र (caritra)}]{\hi{चरित्र} (caritra) — character; moral nature}

\label{lesson:HI-C475-caritra}



\begin{quote}
Before the new one: say the Hindi for mad, then the Hindi for to go astray.
\end{quote}

\subsection*{You'll want to know: \hi{चरित्र}}
\textbf{\hi{चरित्र}} — \emph{caritra} — \textquotedblleft{}character; moral nature\textquotedblright{}.

\begin{quote}
\hi{मीरा} \hi{के} \hi{दादा} \hi{का} \hi{चरित्र} \hi{बहुत} \hi{साफ़} \hi{था।}
\end{quote}

\begin{quote}
\emph{Mīrā ke dādā kā caritra bahut sāf thā.}
\end{quote}

\begin{quote}
Mira's grandfather was a man of spotless character.
\end{quote}

Masculine.

\begin{cousinweb}
\textbf{\hi{चरित्र}} (\emph{caritra}) is a tatsama, borrowed from Sanskrit \emph{caritra}, \textquotedblleft{}acting, behaving, behaviour, deeds\textquotedblright{} (Monier-Williams), from \emph{car}, \textquotedblleft{}to move, to act\textquotedblright{}.
\end{cousinweb}

\begin{grammarlens}[title={What you've built}]
One more word for describing what you lived and felt.
\end{grammarlens}

\subsection*{Your turn}
\begin{itemize}
\raggedright
  \item \emph{Say it:} \emph{caritra}
  \item \emph{Say it:} \emph{caritra}, once more
  \item \emph{Say it:} \emph{Mīrā ke dādā kā caritra bahut sāf thā.} — Mira's grandfather was a man of spotless character.
  \item \emph{Recall:} say the Hindi for mad, then the Hindi for to go astray, then say \emph{caritra} again
\end{itemize}

\begin{tcolorbox}[breakable,colback=teal!4,colframe=teal!35!black,title={Before you move on}]
What does \hi{चरित्र} mean? (\textquotedblleft{}Character; moral nature\textquotedblright{}.) Say it once more.
\end{tcolorbox}
\section[\texorpdfstring{\hi{आचरण} (ācraṇ)}{आचरण (ācraṇ)}]{\hi{आचरण} (ācraṇ) — moral conduct; behaviour}

\label{lesson:HI-C475-acran}



\begin{quote}
Before the new one: say the Hindi for to go astray, then the Hindi for character.
\end{quote}

\subsection*{You'll want to know: \hi{आचरण}}
\textbf{\hi{आचरण}} — \emph{ācraṇ} — \textquotedblleft{}moral conduct; behaviour\textquotedblright{}.

\begin{quote}
\hi{अध्यापक} \hi{बच्चों} \hi{के} \hi{आचरण} \hi{पर} \hi{ध्यान} \hi{देते} \hi{हैं।}
\end{quote}

\begin{quote}
\emph{adhyāpak baccoṁ ke ācraṇ par dhyān dete haiṁ.}
\end{quote}

\begin{quote}
Teachers pay attention to how the children behave.
\end{quote}

Masculine.

\begin{cousinweb}
\textbf{\hi{आचरण}} (\emph{ācraṇ}) is a tatsama, borrowed from Sanskrit \emph{ācaraṇa}, \textquotedblleft{}conduct, behaviour\textquotedblright{} (Monier-Williams), from the same \emph{car} as \emph{caritra}.
\end{cousinweb}

\begin{grammarlens}[title={What you've built}]
One more word for describing what you lived and felt.
\end{grammarlens}

\subsection*{Your turn}
\begin{itemize}
\raggedright
  \item \emph{Say it:} \emph{ācraṇ}
  \item \emph{Say it:} \emph{ācraṇ}, once more
  \item \emph{Say it:} \emph{adhyāpak baccoṁ ke ācraṇ par dhyān dete haiṁ.} — Teachers pay attention to how the children behave.
  \item \emph{Recall:} say the Hindi for to go astray, then the Hindi for character, then say \emph{ācraṇ} again
\end{itemize}

\begin{tcolorbox}[breakable,colback=teal!4,colframe=teal!35!black,title={Before you move on}]
What does \hi{आचरण} mean? (\textquotedblleft{}Moral conduct; behaviour\textquotedblright{}.) Say it once more.
\end{tcolorbox}
\section[\texorpdfstring{\hi{मिज़ाज} (mizāj)}{मिज़ाज (mizāj)}]{\hi{मिज़ाज} (mizāj) — mood; temperament}

\label{lesson:HI-C475-mizaj}



\begin{quote}
Before the new one: say the Hindi for character, then the Hindi for moral conduct.
\end{quote}

\subsection*{You'll want to know: \hi{मिज़ाज}}
\textbf{\hi{मिज़ाज}} — \emph{mizāj} — \textquotedblleft{}mood; temperament\textquotedblright{}.

\begin{quote}
\hi{आज} \hi{दादा} \hi{का} \hi{मिज़ाज} \hi{अच्छा} \hi{नहीं} \hi{है।}
\end{quote}

\begin{quote}
\emph{āj dādā kā mizāj acchā nahīṃ hai.}
\end{quote}

\begin{quote}
Grandfather isn't in a good mood today.
\end{quote}

Masculine.

\begin{cousinweb}
\textbf{\hi{मिज़ाज}} (\emph{mizāj}) is usually traced to Arabic \emph{mizāj}, \textquotedblleft{}mixture\textquotedblright{}, and so \textquotedblleft{}temperament\textquotedblright{}, the mix of humours that old medicine said made a person's nature; it reached Hindi through Persian.
\end{cousinweb}

\begin{grammarlens}[title={What you've built}]
One more word for describing what you lived and felt.
\end{grammarlens}

\subsection*{Your turn}
\begin{itemize}
\raggedright
  \item \emph{Say it:} \emph{mizāj}
  \item \emph{Say it:} \emph{mizāj}, once more
  \item \emph{Say it:} \emph{āj dādā kā mizāj acchā nahīṃ hai.} — Grandfather isn't in a good mood today.
  \item \emph{Recall:} say the Hindi for character, then the Hindi for moral conduct, then say \emph{mizāj} again
\end{itemize}

\begin{tcolorbox}[breakable,colback=teal!4,colframe=teal!35!black,title={Before you move on}]
What does \hi{मिज़ाज} mean? (\textquotedblleft{}Mood; temperament\textquotedblright{}.) Say it once more.
\end{tcolorbox}
\section[\texorpdfstring{\hi{गुण} (guṇ)}{गुण (guṇ)}]{\hi{गुण} (guṇ) — a virtue; a good quality}

\label{lesson:HI-C475-gun}



\begin{quote}
Before the new one: say the Hindi for moral conduct, then the Hindi for mood.
\end{quote}

\subsection*{You'll want to know: \hi{गुण}}
\textbf{\hi{गुण}} — \emph{guṇ} — \textquotedblleft{}a virtue; a good quality\textquotedblright{}.

\begin{quote}
\hi{धैर्य} \hi{अरुण} \hi{का} \hi{सबसे} \hi{बड़ा} \hi{गुण} \hi{है।}
\end{quote}

\begin{quote}
\emph{dhairya Aruṇ kā sabse baṛā guṇ hai.}
\end{quote}

\begin{quote}
Patience is Arun's greatest virtue.
\end{quote}

Masculine.

\begin{cousinweb}
\textbf{\hi{गुण}} (\emph{guṇ}) is a tatsama, borrowed from Sanskrit \emph{guṇa}, first \textquotedblleft{}a single strand of a cord\textquotedblright{} and then \textquotedblleft{}a quality\textquotedblright{} and \textquotedblleft{}good quality, virtue, merit\textquotedblright{} (Monier-Williams).
\end{cousinweb}

\begin{grammarlens}[title={What you've built}]
One more word for describing what you lived and felt.
\end{grammarlens}

\subsection*{Your turn}
\begin{itemize}
\raggedright
  \item \emph{Say it:} \emph{guṇ}
  \item \emph{Say it:} \emph{guṇ}, once more
  \item \emph{Say it:} \emph{dhairya Aruṇ kā sabse baṛā guṇ hai.} — Patience is Arun's greatest virtue.
  \item \emph{Recall:} say the Hindi for moral conduct, then the Hindi for mood, then say \emph{guṇ} again
\end{itemize}

\begin{tcolorbox}[breakable,colback=teal!4,colframe=teal!35!black,title={Before you move on}]
What does \hi{गुण} mean? (\textquotedblleft{}A virtue; a good quality\textquotedblright{}.) Say it once more.
\end{tcolorbox}
\section[\texorpdfstring{\hi{प्रवृत्ति} (pravṛtti)}{प्रवृत्ति (pravṛtti)}]{\hi{प्रवृत्ति} (pravṛtti) — a tendency; a disposition}

\label{lesson:HI-C475-pravrtti}



\begin{quote}
Before the new one: say the Hindi for mood, then the Hindi for a virtue.
\end{quote}

\subsection*{You'll want to know: \hi{प्रवृत्ति}}
\textbf{\hi{प्रवृत्ति}} — \emph{pravṛtti} — \textquotedblleft{}a tendency; a disposition\textquotedblright{}.

\begin{quote}
\hi{बच्चों} \hi{में} \hi{नई} \hi{चीज़} \hi{सीखने} \hi{की} \hi{प्रवृत्ति} \hi{होती} \hi{है।}
\end{quote}

\begin{quote}
\emph{baccoṁ meṁ naī cīz sīkhne kī pravṛtti hotī hai.}
\end{quote}

\begin{quote}
Children have a natural tendency to learn new things.
\end{quote}

Feminine.

\begin{cousinweb}
\textbf{\hi{प्रवृत्ति}} (\emph{pravṛtti}) is a tatsama, borrowed from Sanskrit \emph{pravṛtti}, first \textquotedblleft{}moving onwards\textquotedblright{} and then \textquotedblleft{}a course or tendency towards, an inclination\textquotedblright{} (Monier-Williams).
\end{cousinweb}

\begin{grammarlens}[title={What you've built}]
One more word for describing what you lived and felt.
\end{grammarlens}

\subsection*{Your turn}
\begin{itemize}
\raggedright
  \item \emph{Say it:} \emph{pravṛtti}
  \item \emph{Say it:} \emph{pravṛtti}, once more
  \item \emph{Say it:} \emph{baccoṁ meṁ naī cīz sīkhne kī pravṛtti hotī hai.} — Children have a natural tendency to learn new things.
  \item \emph{Recall:} say the Hindi for mood, then the Hindi for a virtue, then say \emph{pravṛtti} again
\end{itemize}

\begin{tcolorbox}[breakable,colback=teal!4,colframe=teal!35!black,title={Before you move on}]
What does \hi{प्रवृत्ति} mean? (\textquotedblleft{}A tendency; a disposition\textquotedblright{}.) Say it once more.
\end{tcolorbox}
\section[(dialogue)]{Looking back: from Patience and Temper to Character and Conduct}

\label{lesson:HI-R475-review-468-475}



\begin{quote}
Say the Hindi for forbearance and impatient.
\end{quote}

\subsection*{Your turn}
Say these aloud:

\begin{itemize}
\raggedright
  \item forbearance, impatient, cranky, to snap at, to snap irritably
  \item miserly, thrifty, a donation, greedy, a beggar
  \item daring, valiant, bold, fearless, a risk
  \item polite, rude, consideration, etiquette, to tease
  \item nimble, sluggish, good-for-nothing, dedicated, alert
  \item talkative, silent, idle chat, to grumble, silence
  \item sense, a fool, naive, mad, to go astray
  \item character, moral conduct, mood, a virtue, a tendency
\end{itemize}

\begin{tcolorbox}[breakable,colback=teal!4,colframe=teal!35!black,title={Before you move on}]
Forbearance? (\textbf{\hi{धैर्य}}, \emph{dhairya}.) Impatient? (\textbf{\hi{अधीर}}, \emph{adhīr}.) Cranky? (\textbf{\hi{झक्की}}, \emph{jhakkī}.) To snap at? (\textbf{\hi{झिड़कना}}, \emph{jhiṛaknā}.) To snap irritably? (\textbf{\hi{झल्लाना}}, \emph{jhallānā}.) Miserly? (\textbf{\hi{कृपण}}, \emph{kṛpaṇ}.) Thrifty? (\textbf{\hi{मितव्ययी}}, \emph{mitvyayī}.) A donation? (\textbf{\hi{दान}}, \emph{dān}.) Greedy? (\textbf{\hi{लोभी}}, \emph{lobhī}.) A beggar? (\textbf{\hi{भिखारी}}, \emph{bhikhārī}.) Daring? (\textbf{\hi{साहस}}, \emph{sāhas}.) Valiant? (\textbf{\hi{वीर}}, \emph{vīr}.) Bold? (\textbf{\hi{दिलेर}}, \emph{diler}.) Fearless? (\textbf{\hi{निर्भीक}}, \emph{nirbhīk}.) A risk? (\textbf{\hi{जोखिम}}, \emph{jokhim}.) Polite? (\textbf{\hi{विनम्र}}, \emph{vinamra}.) Rude? (\textbf{\hi{बदतमीज़}}, \emph{badtamīz}.) Consideration? (\textbf{\hi{लिहाज़}}, \emph{lihāz}.) Etiquette? (\textbf{\hi{शिष्टाचार}}, \emph{śiṣṭācār}.) To tease? (\textbf{\hi{चिढ़ाना}}, \emph{ciṛhānā}.) Nimble? (\textbf{\hi{फुर्तीला}}, \emph{phurtīlā}.) Sluggish? (\textbf{\hi{सुस्त}}, \emph{sust}.) Good-for-nothing? (\textbf{\hi{निकम्मा}}, \emph{nikammā}.) Dedicated? (\textbf{\hi{कर्मठ}}, \emph{karmaṭh}.) Alert? (\textbf{\hi{मुस्तैद}}, \emph{mustaid}.) Talkative? (\textbf{\hi{बातूनी}}, \emph{bātūnī}.) Silent? (\textbf{\hi{चुप}}, \emph{cup}.) Idle chat? (\textbf{\hi{गप}}, \emph{gap}.) To grumble? (\textbf{\hi{बड़बड़ाना}}, \emph{baṛbaṛānā}.) Silence? (\textbf{\hi{मौन}}, \emph{maun}.) Sense? (\textbf{\hi{अक़्ल}}, \emph{aql}.) A fool? (\textbf{\hi{बेवकूफ़}}, \emph{bevkūf}.) Naive? (\textbf{\hi{नादान}}, \emph{nādān}.) Mad? (\textbf{\hi{पागल}}, \emph{pāgal}.) To go astray? (\textbf{\hi{बहकना}}, \emph{bahaknā}.) Character? (\textbf{\hi{चरित्र}}, \emph{caritra}.) Moral conduct? (\textbf{\hi{आचरण}}, \emph{ācraṇ}.) Mood? (\textbf{\hi{मिज़ाज}}, \emph{mizāj}.) A virtue? (\textbf{\hi{गुण}}, \emph{guṇ}.) A tendency? (\textbf{\hi{प्रवृत्ति}}, \emph{pravṛtti}.)
\end{tcolorbox}
